\NeedsTeXFormat{LaTeX2e}

\documentclass[portrait]{a0poster}

\usepackage{epsfig}
\usepackage{color}
\usepackage{transparent}
\usepackage{wallpaper}
\usepackage{amsmath}
\usepackage{amssymb}
\usepackage{calc}
\usepackage{bm}
\usepackage[utf8]{inputenc}
\usepackage{textcomp}
\usepackage{mathptmx}
\usepackage[scaled=.90]{helvet}
\usepackage{courier}
\usepackage{listings}

\hyphenation{re-so-na-tors na-no-tubes}

\pagestyle{empty}
\TileWallPaper{10.234cm}{11.594cm}{images/dokutitle}

\setlength{\parskip}{5pt plus 2pt minus 1pt}
\setlength{\parindent}{0pt}

\frenchspacing
\sloppy

\renewcommand{\familydefault}{\sfdefault}

% some color definitions
\definecolor{darkblue}{rgb}{0, 0, 0.5}
\definecolor{lightblue}{rgb}{0, 0, 0.8}
\definecolor{lightgray}{gray}{0.75}

\newcommand{\blue}{\color[rgb]{0.2,0.2,1}}
\newcommand{\lightblue}{\color{lightblue}}
\newcommand{\red}{\color{red}}
\newcommand{\darkgreen}{\color[rgb]{0,0.65,0}}

% Delft blauw 
\definecolor{background}{rgb}{1,1,1}
\definecolor{text}{rgb}{0, 0, 0.3}
\definecolor{heading}{rgb}{0, 0, 0.8}

\definecolor{invheading}{rgb}{0.8,0.8,1}

\pagecolor{black}

\lstset{
  showstringspaces=false,
  basicstyle=\ttfamily,
  identifierstyle=\color[rgb]{0,0,0.5},
  stringstyle=\color[rgb]{0.5,0,0.5}\textbf,
  commentstyle=\color[rgb]{0.9,0,0}\textbf
}




% large heading in text
\newcommand{\heading}[1]{
  {\color{heading}\boldmath\textbf{\huge #1}}\\[\medskipamount]
}

% small heading in text
\newcommand{\smallheading}[1]{
  {\color{heading}\textbf{\Large #1}}\\
}

% very small heading in text
\newcommand{\verysmallheading}[1]{
  {\color{heading}\textbf{\large #1}}\\
}

% white bar across the column
\newcommand{\whitebar}{%
\vspace*{0.5cm}
\hspace*{-1.5cm}
\fcolorbox{white}{white}{\rule{30cm}{-0.7cm}}

\vspace*{1cm}
}

% white background for plots etc.
\newcommand{\whitebg}[1]{
 {
  \setlength{\fboxsep}{5mm}
  \setlength{\fboxrule}{1pt}
  \fcolorbox{black}{white}{\color{black}{#1}}
 }
}

% place a figure
\newcommand{\placefigure}[1]{ 
 {
  \vspace{-0.4cm}
  \begin{center}
  \hspace*{-1cm}\begin{minipage}{0.96\textwidth}
   \whitebg{
     \epsfig{#1,width=0.96\textwidth}
   }
   \vspace*{-1cm}
   {\\}
  \end{minipage}
  \end{center} 
  \vspace{-1cm}
 }
}

% place a listing
\newcommand{\placelisting}[1]{ 
 {
  \vspace{-0.2cm}
  \begin{center}
  \hspace*{-1cm}\begin{minipage}{0.96\textwidth}
   \footnotesize \renewcommand{\baselinestretch}{0.85}
   \whitebg{
   \parbox{0.97\textwidth}{
    \lstinputlisting[language=Perl]{#1}
   }
   }
   \renewcommand{\baselinestretch}{1}
   \vspace*{-1cm}
   {\\}
  \end{minipage}
  \end{center} 
  \vspace{-1cm}
 }
}

% place a listing
\newcommand{\placesmalllisting}[1]{ 
 {
  \vspace{-0.2cm}
  \begin{center}
  \hspace*{-1cm}\begin{minipage}{0.96\textwidth}
   \scriptsize \renewcommand{\baselinestretch}{0.8}
   \whitebg{
   \parbox{0.97\textwidth}{
    \lstinputlisting[language=Perl]{#1}
   }
   }
   \renewcommand{\baselinestretch}{1}
   \vspace*{-1cm}
   {\\}
  \end{minipage}
  \end{center} 
  \vspace{-1cm}
 }
}

% place a text
\newcommand{\placetext}[1]{ 
 {
  \vspace{-0.2cm}
  \begin{center}
  \hspace*{-1cm}\begin{minipage}{0.96\textwidth}
   \footnotesize \renewcommand{\baselinestretch}{0.85}
   \whitebg{
   \parbox{0.97\textwidth}{
    \lstinputlisting{#1}
   }
   }
   \renewcommand{\baselinestretch}{1}
   \vspace*{-1cm}
   {\\}
  \end{minipage}
  \end{center} 
 }
}




\begin{document}

\color{text}
\Large
%
\newlength{\seplength}
\newlength{\headerheight}
\newlength{\columnheight}
\newlength{\columnheighta}
\newlength{\columnheightb}
%
\setlength{\columnheight}{104cm}           % value for portrait 3-col mode
\setlength{\columnheighta}{\columnheight-6cm}
\setlength{\columnheightb}{\columnheight+6.5cm}
\setlength{\columnwidth}{0.31\textwidth}   % value for portrait 3-col mode
%
\setlength{\fboxsep}{10mm}
\setlength{\fboxrule}{0mm}
%
%
%
% the page header
%
\hspace*{-1cm}\begin{minipage}[t][][t]{\textwidth-2\fboxsep-8\fboxrule}
\color{invheading}
\begin{center}
{
\VERYHuge \vspace*{-0.6cm}
\textsf{\textbf{
Lab::Measurement -- measurement control with Perl
}}
}
\\[\baselineskip]
\begin{tabular}{ccc}
\parbox{5cm}{\vspace*{-1cm}
\begin{center}
\epsfig{file=images/logo-ur,height=4.5cm}
\end{center}
}
&
\parbox{70cm}{
\begin{center} 
{
\vspace*{-0.5cm}
\huge 
S. Reinhardt$^1$, C. E. Lane$^2$, C. Butschkow$^1$, A. Iankilevitch$^1$, A. 
Dirnaichner$^1$, 
and A. K. Hüttel$^1$
}
\\[\medskipamount]
\it 
\Large
$^1$Institute for Experimental and Applied
Physics, University of Regensburg, 93040 Regensburg, Germany\\
$^2$Department of Physics, Drexel University, 3141 Chestnut Street, 
Philadelphia, PA 19104, USA\\
\end{center}
\vspace*{0.5cm}
}
&
\parbox{5cm}{\vspace*{-1cm}
\begin{center}
\epsfig{file=images/logo-sfb689,height=4.5cm}
\end{center}
}
\end{tabular}
\end{center}
\end{minipage}
\vspace*{9mm}





%
% begin of first row first column
\fcolorbox{white}{background}{
\begin{minipage}[t][\columnheighta-2\fboxsep-2\fboxrule][t]
  {\columnwidth-2\fboxsep-2\fboxrule} \rule{0pt}{0pt}\\
  \begin{minipage}{\textwidth}
%
%

\vspace*{0.9cm}
\heading{Flexible measurement needed?!}

\vspace*{-1cm}
\begin{itemize}
\item
Tired of following your wires in square meters of LabVIEW diagrams?
\item
Tired of clumsy string handling and low-level driver functions
in your looong C program?
\item
Use a text processing language to manage your measurement! Use Perl!
\end{itemize}
\vspace*{0.5cm}

\placelisting{srs_read.pl}

\vspace*{2cm}
\vspace*{0.9cm}
\heading{Currently supported hardware}

\vspace*{-3mm}
\placefigure{file=images/hardware}
Hardware driver backends: 
\begin{itemize}
\item 
  NI-VISA (MS Windows) and all hardware supported by it
\item
  LinuxGPIB and all hardware supported by it
\item
  TCP connection, generic network socket (both Linux and Windows)
\item
  USB-TMC lightweight driver (Linux/Unix, libusb)
\item
  VXI-11 lightweight driver (Linux/Unix, libtirpc)
\item
  Zurich Instruments LabOne API (both Linux and Windows)
\end{itemize}

\vspace*{1cm}
Growing number of high-level drivers (more are {\it very} easy to add):
{\large
\begin{itemize}
\item Multimeters: HP / Agilent / Keysight
\item DC sources: Yokogawa / Keithley
\item Lock-in amplifiers: Stanford Research / Signal Recovery / Zurich Instruments
\item Temperature controllers: Lakeshore / Oxford Instruments
\item RF / microwave sources, spectrum analyzers, VNAs from Rohde \& Schwarz 
\item and many more...
\end{itemize}
}
\vspace*{1cm}
\heading{Key facts}
\vspace*{-2cm}
\begin{itemize}
 \item
 Open source / free software\\[-5cm]
 \item 
 http://www.labmeasurement.de/
 \hfill\epsfig{file=images/qrcode-labmeasurement,width=5cm}
 \item
 License: same as Perl (GPL-1+ or Artistic)
 \item
 Releases on CPAN, development on Github
 \item
 Contributors and cooperations welcome!
\end{itemize}

%
%
%
\end{minipage}
%
\end{minipage}}
% end of first row first column
%
%
%
%
% begin of first row second column
\fcolorbox{white}{background}{
\begin{minipage}[t][\columnheighta-2\fboxsep-2\fboxrule][t]
  {\columnwidth-2\fboxsep-2\fboxrule} \rule{0pt}{0pt}\\
  \begin{minipage}{\textwidth}
%
%
%
%

\vspace*{0.9cm}
\heading{Real world measurement}

\vspace*{-1cm}
\begin{itemize}
 \item 
Ferromagnetic resonance measurement with Vector Network Analyzer (VNA)
 \item
Outer loop: Continuous sweep of magnetic field with OI Mercury iPS
 \item 
Inner loop: VNA transmission measurement at multiple microwave frequencies
\end{itemize}

\vspace*{0.5cm}
%\placelisting{VNA-FMR.pl}
 \placesmalllisting{VNA-FMR.pl}
\vspace*{3cm}

\heading{Output files}
\vspace*{-1cm}
\placelisting{ls.txt}
\vspace*{2cm}
\begin{itemize}
\item \emph{META.yml:} Various metadata (host, user, date, L::M version,
  commandline arguments)
\item \emph{data.dat:} Measured data, tab-separated gnuplot format
\item \emph{data.dat.plot.png:} Live plot at end of measurement as png image
\item \emph{magnetic-resonance.pl:} Archival copy of the measurement script

\end{itemize}

\heading{Advanced Sweep features}
\vspace*{-2cm}
\begin{itemize}
\item Higher dimensional sweeps, e.g. 3D sweep: Create one 2D datafile for each
  step of the outer sweep
\item Log arrays and matrices of data (PDLs). Useful with spectrum analyzers,
  VNAs, and oscilloscopes
\item Extensive support for live plots via gnuplot: Line plots (2D data) and color maps (3D
  data)
\item Customize live plots: access to all of gnuplot's plot and curve options
  --- both for live plots (qt, x11, wxt) and hardcopies 
(png, jpeg, svg, tikz, pdfcairo, ...) via PDL::Graphics::Gnuplot

\end{itemize}



% 
% 
% 
% 
\end{minipage}
\end{minipage}}
% end of second column
%
%
%
%
%
%begin third column
\fcolorbox{white}{background}{
\begin{minipage}[t][\columnheighta-2\fboxsep-2\fboxrule][t]
  {\columnwidth-2\fboxsep-2\fboxrule} \rule{0pt}{0pt}\\
  \begin{minipage}{\textwidth}
%
%
%
%
\vspace*{0.9cm}

\heading{Instrument/Connection\\abstraction layers}

\vspace*{-3mm}
\placefigure{file=images/structure}

\begin{itemize}
\item
Modular structure. Easy to extend with new instrument drivers and connection
types.
\item
Abstract IO layer, makes instrument drivers independent of hardware backends.
\end{itemize}

\vspace*{2cm}
\heading{Recent improvements}
\vspace*{-2cm}
\begin{itemize}
\item {\lightblue New drivers}: OI Mercury magnet controller and level meter
  sensor; Zurich Instruments MFLI lock-in and MFIA impedance analyzer
\item {\lightblue New connection backends}
  \begin{itemize}
  \item USB-TMC, VXI-11: Lightweight user space drivers for Linux/Unix.
    Provide {\lightblue fast and cheap alternatives to GPIB} on Linux. High
    reliablitiy as based on technically mature system libraries (libusb, sunrpc)
  \item Initial support for LabOne API of Zurich Instruments: Provide Perl
    bindings to proprietary vendor library. Instruments are accessed via
    Ethernet or USB
  \end{itemize}
\item {\lightblue New high-level Sweep framework:} Modern Perl
  reimplementation of previous Lab::XPRESS layer
  \vspace*{0.5cm}
  \placefigure{file=images/moose-sweep}
  \begin{itemize}
  \item Separate classes for sweeps, datafiles, datafolder, and plots
  \item Most operational details of the sweeps are implemented in subclasses of
    the Lab::Moose::Sweep base class
    \item High modularity: \emph{Very} easy to extend with new functionality
  \end{itemize}
\end{itemize}

\vspace*{2cm}
\heading{Outlook}
\vspace*{-2cm}

\begin{itemize}
\item Implement more sweep types with the new framework, e.g. temperature sweep
\item {\lightblue Improve on tutorials:}
  \begin{itemize}
  \item Entry-level for new users: Basics about connections, instruments,
    and sweeps
  \item Developer docs: How to write and test new drivers; use Moose features
    efficiently
  \end{itemize}
  
\end{itemize}


\end{minipage}

\end{minipage}}  % end of first row third column
%
%
%

\vspace*{0.5cm}
\begin{center}\color{invheading}\large
\hspace*{-4cm}We gratefully acknowlegde funding by the DFG via the Emmy Noether grant
Hu1808/1-1 ``Carbon nanotubes as electronical and nano-electromechanical hybrid systems in
the quantum limit'',
\\ \hspace*{-4cm}
the collaborative research centre SFB 689 ``Spin phenomena in reduced dimensions'', and
the graduate research school GRK 1570  ``Electronic properties of
carbon based nanostructures''.
\end{center}

\end{document}
